\textbf{1. Kiến thức cần nhớ}
\begin{itemize}
\item \textbf{Khoảng biến thiên ($R$):} Là hiệu giữa giá trị lớn nhất và giá trị nhỏ nhất của mẫu số liệu:
\[ \boxed{R = x_{\max} - x_{\min}} \]
\textit{Ý nghĩa:} Đo độ phân tán của toàn bộ mẫu số liệu, nhưng dễ bị ảnh hưởng bởi các giá trị bất thường.
\item \textbf{Khoảng tứ phân vị ($\Delta_Q$):} Là hiệu giữa tứ phân vị thứ ba và tứ phân vị thứ nhất:
\[ \boxed{\Delta_Q = Q_3 - Q_1} \]
\textit{Ý nghĩa:} Đo độ phân tán của $50\%$ số liệu chính giữa, không bị ảnh hưởng bởi các giá trị bất thường.
\item \textbf{Giá trị bất thường (ngoại lệ):} Giá trị $x$ trong mẫu số liệu được coi là giá trị bất thường nếu:
\[ \boxed{x > Q_3 + 1{,}5\Delta_Q \quad \text{hoặc} \quad x < Q_1 - 1{,}5\Delta_Q} \]
\item \textbf{Phương sai ($s^2$):} Cho mẫu số liệu $x_1, x_2, \ldots, x_n$ có số trung bình $\overline{x}$:
\[ \boxed{s^2 = \dfrac{1}{n}\left[(x_1 - \overline{x})^2 + (x_2 - \overline{x})^2 + \cdots + (x_n - \overline{x})^2\right] = \dfrac{x_1^2 + x_2^2 + \cdots + x_n^2}{n} - \overline{x}^2} \]
Nếu mẫu cho dưới dạng bảng tần số (giá trị $x_i$ có tần số $n_i$, cỡ mẫu $n = \sum n_i$):
\[ s^2 = \dfrac{1}{n}\sum n_i (x_i - \overline{x})^2 = \dfrac{1}{n}\sum n_i x_i^2 - \overline{x}^2 \]
\item \textbf{Độ lệch chuẩn ($s$):} Là căn bậc hai của phương sai: $\boxed{s = \sqrt{s^2}}$.
\textit{Ý nghĩa:} Phương sai và độ lệch chuẩn đo mức độ phân tán của các số liệu quanh số trung bình $\overline{x}$; giá trị càng nhỏ thì mẫu số liệu càng đồng đều (ổn định). Độ lệch chuẩn có cùng đơn vị với mẫu số liệu.
\end{itemize}

\textbf{2. Phương pháp giải các dạng}

%% DANG 0D6?4-1
\textbf{Dạng 1. Câu hỏi lý thuyết}
\begin{enumerate}
\item Xác định số đặc trưng đo mức độ phân tán được đề cập ($R$, $\Delta_Q$, $s^2$ hay $s$).
\item Đối chiếu định nghĩa, công thức, đơn vị đo và ý nghĩa thực tiễn để chọn đáp án đúng.
\end{enumerate}
Lưu ý: Khi mẫu có giá trị bất thường, nên dùng $\Delta_Q$ thay cho $R$, $s^2$, $s$; hai mẫu có số trung bình bằng nhau thì mẫu nào có $s$ nhỏ hơn sẽ ổn định hơn.

%% DANG 0D6?4-2
\textbf{Dạng 2. Khoảng biến thiên, khoảng tứ phân vị}
\begin{enumerate}
\item Sắp xếp mẫu số liệu theo thứ tự không giảm $x_1 \le x_2 \le \cdots \le x_n$.
\item Tính khoảng biến thiên $R = x_n - x_1$.
\item Tìm các tứ phân vị $Q_1, Q_2, Q_3$ rồi tính khoảng tứ phân vị $\Delta_Q = Q_3 - Q_1$.
\end{enumerate}
Lưu ý: Với bảng tần số, $x_{\min}$ và $x_{\max}$ là giá trị nhỏ nhất và lớn nhất có tần số dương, không phải tần số nhỏ nhất/lớn nhất.

%% DANG 0D6?4-3
\textbf{Dạng 3. Giá trị bất thường của mẫu số liệu}
\begin{enumerate}
\item Sắp xếp mẫu số liệu không giảm, tìm $Q_1, Q_3$ và tính $\Delta_Q = Q_3 - Q_1$.
\item Tính hai ngưỡng $A = Q_1 - 1{,}5\Delta_Q$ và $B = Q_3 + 1{,}5\Delta_Q$.
\item Kiểm tra các giá trị trong mẫu: mọi giá trị nằm ngoài đoạn $[A; B]$ đều là giá trị bất thường.
\end{enumerate}
Lưu ý: Giá trị bằng đúng $Q_1 - 1{,}5\Delta_Q$ hoặc $Q_3 + 1{,}5\Delta_Q$ không phải là giá trị bất thường (bất đẳng thức ngặt $<$ hoặc $>$).

%% DANG 0D6?4-4
\textbf{Dạng 4. Phương sai, độ lệch chuẩn của mẫu số liệu}
\begin{enumerate}
\item Tính số trung bình $\overline{x}$ của mẫu số liệu.
\item Tính phương sai $s^2$ bằng định nghĩa hoặc công thức rút gọn $s^2 = \dfrac{1}{n}\sum x_i^2 - \overline{x}^2$ (hoặc bấm máy tính cầm tay).
\item Lấy căn bậc hai để tìm độ lệch chuẩn $s = \sqrt{s^2}$.
\end{enumerate}
Lưu ý: Khi dùng máy tính cầm tay theo sách Kết nối tri thức, phương sai là $\sigma_x^2$ và độ lệch chuẩn là $\sigma_x$ (không nhầm với $s_x^2$ và $s_x$ chia cho $n-1$).
